\section{Experimental design}

The evaluation was designed to answer one methodological question: can the
framework distinguish useful prediction and reproducibility from evidence that
supports a bounded mechanistic claim? Four experimental regimes were selected
because they expose different parts of this question. Their outputs are not
pooled into a single performance table. Instead, each regime tests a declared
failure mode or evidence block.

\subsection{General protocol and reproducibility controls}

Every experiment is executed from a versioned configuration or command and
writes a structured result artifact. Cohort size, data identifiers, thresholds,
model settings, random seeds, and software revision are stored with the result.
Derived publication tables are generated from these artifacts. No raw public
dataset is redistributed, and acquisition scripts preserve source URLs or API
queries so that authorized users can recreate each local cache.

Thresholds define minimum evidence for the target claim rather than tunable
objectives. Sensitivity analysis perturbs relevant thresholds and checks whether
the qualitative conclusion changes. Negative outcomes are retained. This is
important because omitting failed topology or directionality tests would make
the audit indistinguishable from selecting only favorable metrics.

For pairwise analyses, multiplicity and dependence receive separate treatment.
False-discovery-rate correction controls the family of endpoint tests. Bootstrap
resampling is performed over units, not over the much larger set of derived
directed pairs, because pairs sharing a neuron are not independent. MICRONS
replication uses non-overlapping query windows that are fixed before inspecting
the hold-out results.

\subsection{Known-truth validation of MIS}

Four synthetic systems isolate the logical components of identifiability. The
\texttt{directed\_truth} case contains reproducible region-specific signals,
an informative topology, and resolvable directed latency. The
\texttt{common\_drive\_nonidentifiable} case is reproducible and predictable
but generated by a common input. The \texttt{topology\_without\_direction}
case preserves region-specific structure without a resolvable temporal order.
The \texttt{direction\_without\_topology\_specificity} case contains timing
information while alternative graphs explain the observations equally well.

For each system, the proposed topology is compared with disconnected and
permuted controls. Reproducibility is measured through cross-session and
split-half correlations and the fraction of instances exceeding their own null
distribution. Topology specificity uses prediction advantages over alternative
graphs. Directed identifiability uses resolvable latency and nonzero lead-lag
fractions. The prespecified expectation is strict: only
\texttt{directed\_truth} may pass all MIS blocks.

This experiment is a unit test of the scientific decision rule, not biological
validation. Its purpose is to expose compensatory behavior. If high
reproducibility allowed a common-drive system to pass, or if topology alone
compensated for absent directionality, the score would be unsuitable for the
remaining experiments.

\subsection{Allen negative mechanistic control}

Allen Visual Behavior Neuropixels is used as a real negative control. The
confirmation analysis contains 20 successful sessions. The anatomical
specificity test contains 20 sessions from 20 unique mice, and the directed
identifiability analysis contains 21 sessions from 19 mice. Population targets
are evaluated across mice and split halves before applying anatomical and
temporal gates.

The topology block compares the Allen-informed model with permuted and
disconnected alternatives. It requires a positive advantage of at least 0.05,
at least 95\% of permutations to be outperformed, and a positive lower
bootstrap bound. The directional block evaluates cross-mouse latency agreement,
split-half latency agreement, resolvable regional pairs, and nonzero lead-lag
structure. These requirements are intentionally stronger than the observation
that an evoked response is repeatable.

Allen is not expected to be a positive connectomic result in this configuration.
The individual synaptic network producing each session is not observed.
Accordingly, the informative outcome is whether MouseBrainBench preserves the
strong reproducibility finding while refusing to infer an anatomical-dynamic
mechanism that the target cannot identify.

\subsection{Sensorium predictive controls}

Static Sensorium is evaluated on seven validation mice and on five mice with
repeated test stimuli. Transparent predictors include a mean-response baseline,
stimulus ridge regression, and stimulus-plus-context ridge regression. The
analysis reports neural reliability, correlation, improvement over the mean
response, and improvement over a scrambled-stimulus control. A topographic test
is evaluated separately because spatial organization is structural evidence but
not directed connectivity.

Dynamic Sensorium is evaluated on five current public mice and five legacy
out-of-distribution mice. The model set includes mean response, summary
features, a temporal filterbank, temporal singular-value decomposition, random
features, and a compact multilayer perceptron. An official Sensorium loader and
model factory are also integrated under a bounded local training budget. That
run is an interoperability control; it is not considered leaderboard-equivalent
or state of the art because it does not reproduce the published training budget.

The Sensorium experiments test whether conventional predictive reporting and
the claim audit lead to different conclusions. Correlation and improvement over
simple baselines can establish predictive utility. They cannot pass topology or
directionality blocks without corresponding anatomical or interventional
evidence. The framework should therefore report positive predictive evidence
and a negative mechanistic decision simultaneously.

\subsection{MICRONS structure-function experiment}

MICRONS provides the main empirical test because functional properties and
synaptic edges can be co-registered locally. Data are queried from the
\texttt{minnie65\_public} CAVE stack using
\texttt{digital\_twin\_properties\_bcm\_coreg\_v4} at materialization
version 1507. Three non-overlapping windows are used: discovery offset 0,
hold-out offset 1000, and hold-out offset 2000. They contain 1000, 992, and 999
usable units, respectively.

The primary endpoint is fixed as readout-location similarity among all directed
non-self pairs. Connected pairs are contrasted with non-connected pairs under
three null constructions: random controls, controls matched by intersomatic
distance, and controls matched by neuronal degree. Distance matching addresses
the fact that nearby neurons are both more likely to connect and more likely to
share functional organization. Degree matching addresses unequal opportunities
for high-degree neurons to appear in connected pairs.

The endpoint is admitted only when the connected-pair advantage remains
positive after matched controls and false-discovery-rate correction in the
discovery cohort and both hold-outs. Stability is assessed with 300 unit-cluster
bootstrap samples per cohort using seed 71. This protocol supports, at most, a
replicated local observational association. No perturbation establishes synaptic
necessity, and the sampled volume cannot support a whole-brain conclusion.
